% Source PDF: source/普通高考/2017/2017全国2文(甘肃,青海,内蒙古,黑龙江,吉林,辽宁,海南,宁夏,新疆,陕西,重庆).pdf
% Page: 1 (the two-page A3 source PDF is vector-only; pdfimages found no embedded bitmap).
% Grid evidence: tmp/redraw_flowchart_2017_national_paper_2_liberal/grid/page1_grid100.png and q10_block_grid50.png.
% Comparison source: tmp/redraw_flowchart_2017_national_paper_2_liberal/crops/q10_orig_final.png,
% cropped from the ungridded page render with a safety border.
% Elements: rounded start/end terminals, input/output parallelograms, initialization/update
% rectangles, K<=6 decision diamond, directional arrows, and the yes/no loop branches.
% Relations: start -> input a -> S=0,K=1 -> decision; yes -> S update -> sign flip -> K increment -> decision;
% no -> output S -> end. The left and right polylines are the two control-flow return/exit branches.
% Directed edges (read directly from the PDF): start->input, input->init, init->test,
% test-(是)->sum, sum->flip, flip->inc, inc->test, test-(否)->output, output->finish.
% solid segments: all control-flow connectors and node borders; dashed segments: none; labels are kept
% inside nodes or beside their branch connectors to avoid overlap with arrows.
\begin{tikzpicture}[x=1cm,y=1cm,baseline,
  flowarrow/.style={-{Latex[length=2.1mm,width=1.5mm]},line width=0.45pt},
  nodebox/.style={draw,line width=0.45pt,inner sep=2pt,minimum height=0.62cm,anchor=center,align=center,
    execute at begin node={\setlength{\parindent}{0pt}\noindent}},
  branchlabel/.style={anchor=center,align=center,inner sep=0pt,
    execute at begin node={\setlength{\parindent}{0pt}\noindent}},
  term/.style={nodebox,rounded corners=2.5pt,minimum width=1.05cm},
  io/.style={nodebox,shape=trapezium,trapezium left angle=70,trapezium right angle=110,minimum width=1.85cm,minimum height=0.68cm},
  decision/.style={nodebox,diamond,aspect=2.75,minimum width=3.00cm,minimum height=1.05cm,inner sep=1pt}
]
  % Node layout follows the top-to-bottom control-flow order in the source diagram.
  \node[term] (start) at (0,0) {开始};
  \node[io] (input) at (0,-1.05) {输入 $a$};
  \node[nodebox,minimum width=2.55cm] (init) at (0,-2.20) {$S=0,K=1$};
  \node[decision] (test) at (0,-3.45) {$K\leq 6$};
  \node[nodebox,minimum width=3.15cm] (sum) at (0,-4.68) {$S=S+a\cdot K$};
  \node[nodebox,minimum width=1.80cm] (flip) at (0,-5.82) {$a=-a$};
  \node[nodebox,minimum width=2.35cm] (inc) at (0,-6.98) {$K=K+1$};
  \node[io] (output) at (0,-8.27) {输出 $S$};
  \node[term] (finish) at (0,-9.35) {结束};

  % Main vertical path and the yes branch.
  \draw[flowarrow] (start.south) -- (input.north);
  \draw[flowarrow] (input.south) -- (init.north);
  \draw[flowarrow] (init.south) -- (test.north);
  \draw[flowarrow] (test.south) -- (sum.north);
  \draw[flowarrow] (sum.south) -- (flip.north);
  \draw[flowarrow] (flip.south) -- (inc.north);
  \draw[flowarrow] (output.south) -- (finish.north);

  % Loop-back from K increment: the horizontal segment enters the shared decision stem.
  \draw (inc.west) -- (-2.25,-6.98) -- (-2.25,-2.65);
  \draw[flowarrow] (-2.25,-2.65) -- (0,-2.65);
  \draw (0,-2.65) -- (test.north);
  % The no branch exits right, runs down, then turns left to the output stem.
  \draw (test.east) -- (2.55,-3.45) -- (2.55,-7.70) -- (0,-7.70);
  \draw[flowarrow] (0,-7.70) -- (output.north);
  \node[branchlabel] at (1.18,-3.02) {否};
  \node[branchlabel] at (0.48,-3.98) {是};
\end{tikzpicture}
